problem:  
Let \( N_{k,l} = \mathrm{SU}(3)/\mathrm{U}(1)_{k,l} \), where  
\[
\mathrm{U}(1)_{k,l}=\left\{\operatorname{diag}(e^{ik\theta}, e^{il\theta}, e^{-i(k+l)\theta}) : \theta\in \mathbb{R}\right\}.
\]  
It is known that certain \(N_{k,l}\) admit invariant nearly parallel \(\mathrm{G}_2\)-structures and that their space of invariant \(\mathrm{G}_2\)-structures can be parameterized by a finite-dimensional family.

**Problem:**  
For fixed integers \((k,l)\), consider the moduli space  
\[
\mathcal{M}_{k,l} = \{ \text{SU(3)-invariant } \mathrm{G}_2\text{-structures on } N_{k,l}\} / \text{SU(3)-equivariant diffeomorphisms and rescaling}.
\]
1. Determine whether the subset \(\mathcal{M}_{k,l}^{\text{np}}\) of nearly parallel invariant \(\mathrm{G}_2\)-structures lies in the closure of the subset \(\mathcal{M}_{k,l}^{\text{cc}}\) of coclosed invariant \(\mathrm{G}_2\)-structures in \(\mathcal{M}_{k,l}\). In other words, can a sequence of coclosed invariant \(\mathrm{G}_2\)-structures approach a nearly parallel one?  

2. If such a limiting phenomenon occurs, study whether the scalar curvature of the associated invariant metrics \(g_\varphi\) satisfies  
\[
\liminf_{\varphi\to\varphi_0} \mathrm{Scal}(g_\varphi) > 0,
\]
for \(\varphi_0 \in \mathcal{M}_{k,l}^{\text{np}}\).  

3. Determine for which pairs \((k,l)\) this boundary relation between coclosed and nearly parallel loci in \(\mathcal{M}_{k,l}\) can exist, and whether it is unique or has multiple branches within the moduli topology induced by the invariant parameter space.

Why is it a "good" problem:  
This problem is sharply formulated, conceptually simple, and yet deeply open. It asks whether the space of invariant \(\mathrm{G}_2\)-structures on an Aloff–Wallach space admits a *limiting interface* between two geometrically distinct types—coclosed and nearly parallel—under the natural homogeneous moduli topology. Resolving this would require understanding both the **representation-theoretic structure** of invariant forms on \(N_{k,l}\) and their **curvature functionals**, bridging special holonomy theory with curvature analysis. It serves as a potential model for how Einstein and nearly parallel geometries might arise as degenerations of other \(\mathrm{G}_2\)-structures, connecting **invariant geometry**, **analysis on homogeneous spaces**, and **moduli theory**. The problem’s statement is elementary in form but its resolution would demand profound insights into the analytical and algebraic structure of Aloff–Wallach geometries—capturing the “narrow entrance, deep content” quality of a good research-level problem.